\SGASection{4}{Roots and vector group schemes}
\SGARefTarget{sga3:xix:section:4}

\par\medskip\noindent
\textbf{\SGAInlineRefTarget{sga3:xix:paragraph:4-1}4.1.}---
Let \(S\) be a scheme and let \(F\) be a finite locally free
\(\mathcal O_S\)-module.  The \(S\)-scheme \(\mathbb W(F)\) is smooth
over \(S\).  Its Lie algebra is canonically isomorphic to \(F\).
Indeed, there is a canonical isomorphism
\[
  \mathbb W(F)
  \xrightarrow{\ \sim\ }
  \operatorname{Lie}(\mathbb W(F)/S)
  =
  \mathbb W\bigl(\operatorname{Lie}(\mathbb W(F)/S)\bigr)
\]
\textup{(\SGARefLink{sga3:ii:corollary:4-4-1}{Exposé II, 4.4.1} and \SGARefLink{sga3:xix:scholium:4-4-2}{4.4.2})}.  We shall always identify
\(F\) with \(\operatorname{Lie}(\mathbb W(F)/S)\).
\SGARefTarget{sga3:xix:paragraph:4:1}

\par\medskip\noindent
\textbf{\textit{\SGARefLink{sga3:xix:lemma:4:2}{Lemma 4.2}.}}---
\textit{Let \(S\) be a scheme and let \(V\) be a vector bundle over
\(S\), smooth over \(S\).  There is then a unique isomorphism of
\(\mathcal O_S\)-modules}
\[
  \exp\colon
  \mathbb W\bigl(\operatorname{Lie}(V/S)\bigr)
  \xrightarrow{\ \sim\ } V
\]
\textit{which induces the identity on Lie algebras.}
\footnote{Editor's note: the proof in the original edition has been
retained; it may also be expanded as follows.

Let \(\mathscr F\) be a quasi-coherent \(\mathcal O_S\)-module and put
\(V=\mathbb V(\mathscr F)\).  Write \(\pi\colon V\to S\) for the
projection and \(\epsilon\colon S\to V\) for the zero section.  Then
\(\Omega^1_{V/S}=\pi^*\mathscr F\), whence
\[
  \tag{1}\label{sga3:xix:paragraph:4-1:equation:1}
  \omega^1_{V/S}
  =
  \epsilon^*\Omega^1_{V/S}
  \simeq
  \epsilon^*\pi^*\mathscr F
  \simeq
  \mathscr F,
\]
and therefore
\(\operatorname{Lie}(V/S)=(\omega^1_{V/S})^\vee
\simeq\mathscr F^\vee\).
If \(V\) is smooth over \(S\), then by \textup{(1)}
\(\mathscr F\) is finite locally free, and hence
\[
  \tag{2}\label{sga3:xix:paragraph:4-1:equation:2}
  V=\mathbb V(\mathscr F)
  \simeq
  \mathbb W(\mathscr F^\vee)
  \simeq
  \mathbb W\bigl(\operatorname{Lie}(V/S)\bigr).
\]
Now let \(U\) be an \(S\)-scheme equipped with a section
\(\tau\colon S\to U\), and suppose that an isomorphism
\(\phi\colon V\xrightarrow{\sim}U\) of pointed \(S\)-schemes has been
given, that is, \(\phi\circ\epsilon=\tau\) \textup{(for example, when
\(U\) is an \(S\)-group)}.  Then \(\phi\) induces an isomorphism
\[
  d\phi\colon
  \mathscr F^\vee
  \xrightarrow{\ \sim\ }
  \operatorname{Lie}(U/S),
\]
and therefore an isomorphism \(\psi\) displayed by
\[
\begin{tikzcd}[
  ampersand replacement=\&,
  column sep=huge,
  row sep=large
]
  \mathbb W(\mathscr F^\vee)
    \arrow[r,equal]
  \&
  \mathbb V(\mathscr F)
    \arrow[r,"\phi","\sim"']
  \&
  U
  \\
  \&
  \mathbb W\bigl(\operatorname{Lie}(U/S)\bigr)
    \arrow[ul,"\sim","\mathbb W(d\phi^{-1})"']
    \arrow[ur,"\sim","\psi"']
  \&
\end{tikzcd}
\]
which identifies \(U\) with
\(\mathbb W(\operatorname{Lie}(U/S))\).  Since the functor
\(\mathbb W\) is fully faithful, there is a unique isomorphism
\[
  \exp\colon
  \mathbb W\bigl(\operatorname{Lie}(U/S)\bigr)
  \xrightarrow{\ \sim\ }U
\]
such that
\(\operatorname{Lie}(\psi^{-1}\circ\exp)=\operatorname{id}\).
In fact, one sees directly that
\(\operatorname{Lie}(\psi)=\operatorname{id}\), so that
\(\exp=\psi\).}
\SGARefTarget{sga3:xix:lemma:4:2}

\SGAProof{Proof}
We have \(V=\mathbb V(\mathscr F)\) for some quasi-coherent
\(\mathcal O_S\)-module \(\mathscr F\).  Since \(V\) is smooth over
\(S\), the module
\(\mathscr F\simeq\omega^1_{V/S}\) is finite locally free, and hence
\[
  \operatorname{Lie}(V/S)
  =
  \mathbb V(\omega^1_{V/S})
  \simeq
  \mathbb W\bigl(\operatorname{Lie}(V/S)\bigr).
\]
Moreover, by Exposé II, loc.\ cit., there is a canonical isomorphism
\[
  V
  \xrightarrow{\ \sim\ }
  \operatorname{Lie}(V/S)
  \simeq
  \mathbb W\bigl(\operatorname{Lie}(V/S)\bigr),
\]
and the uniqueness of \(\exp\) follows at once because
\(\mathbb W\) is fully faithful.
\hfill\(\square\)
